\section{The Ark Goes Before Them}

Before the Ark moves, it disappears beneath veils.  Aaron and his sons enter
first, take down the screening curtain, cover the Ark, add protective coverings,
and set its poles in place.  Only then may the Kohathites approach their burden.
They bear the holy things, but they must neither touch them nor press in to gaze
upon them, lest they die (Num 4:5--20).  Other Levitical families receive carts
for their work.  The Kohathites receive none, because the sanctuary entrusted
to them is borne on their shoulders (Num 7:3--9).

These commands give the journey its grammar.  The Ark does not move like
ordinary cargo, and the bearers do not improvise their office.  Holiness creates
an ordered nearness: the priests prepare, the appointed men bear, the poles
carry the weight, and the people follow the Presence God has given.

\subsection{The three-day vanguard}

On the twentieth day of the second month of the second year, the cloud lifts
from the tabernacle and Israel leaves Sinai.  Standards rise; the tribes move in
their appointed order; the dismantled sanctuary travels in the midst of the
camp.  Then Numbers draws the Ark itself into the foreground.  For a three-day
journey the Ark of the Covenant goes before the people to seek a resting place.
When it sets out Moses summons the Lord to arise and scatter his enemies; when
it rests he asks the Lord to return to the myriads of Israel
(Num 10:11--36).

The passage speaks of this three-day advance.  It does not say that the Ark
always occupied the same place in every march, and Numbers 33 lists Israel's
stations rather than an itemized Ark itinerary.  The restraint matters because
the image is already strong enough: before Israel knows its rest, the covenant
sign goes ahead.

At Kadesh the reverse scene is devastating.  After refusing the promised land,
the people decide to advance when the Lord has not sent them.  Moses warns
them, but they climb toward the hill country.  Moses and the Ark remain in the
camp.  The presumptuous assault ends in flight and defeat (Num 13:26;
14:39--45).  The Ark is no mechanism of success.  It can lead obedience; it
cannot be dragged behind repentance that refuses to listen.

\subsection{The river stands open}

At Shittim, on the edge of the land, officers pass through the camp with a new
command: when the people see the priests bearing the Ark, they are to follow,
while keeping the appointed distance.  They have not passed this way before
(Josh 3:1--6).  The words make the crossing more than a maneuver.  Israel must
learn the road by following.

The priests lift the Ark and walk toward the Jordan.  It is harvest time, when
the river is full.  Their feet reach the water; the flow is cut off upstream;
and the bearers stand on dry ground in the middle of the riverbed while all
Israel crosses.  Twelve stones are taken for a memorial.  Only after the whole
command is complete do the priests come up; when their feet reach the bank, the
river returns to its course (Josh 3:7--17; 4:1--18).  The text says that the
priests' feet enter the water.  It does not make the chest itself touch the
river.

Israel camps at Gilgal and sets up the memorial stones there
(Josh 4:19--24).  The Ark has crossed with the people, but Scripture does not
separately narrate its installation at Gilgal.  The named camp is therefore a
strong narrative setting, not a license to supply an unrecovered shrine or
road.

\subsection{Silence, trumpets, and a fallen wall}

Jericho closes its gates.  Israel does not answer first with siege craft.
Armed men go before seven priests carrying seven trumpets; the Ark follows;
the rear guard comes behind.  For six days they circle the city once.  On the
seventh they circle seven times.  The trumpets sound, Joshua gives the word,
the people shout, and the wall falls (Josh 6:1--20).  The same Ark that stood
still in the Jordan now moves through a liturgical circuit.  The action is
God's gift under God's command, not a ritual formula Israel has mastered.

The next campaign proves the point.  Israel is defeated after Achan's covenant
violation, and Joshua falls before the Ark until evening (Josh 7:1--9).  The
narrative does not locate the Ark at Ai; Joshua's place of lament is unnamed.
Presence exposes sin as well as overcoming enemies.

After judgment and renewed obedience, Joshua builds an altar on Mount Ebal and
reads the law before the assembled people.  Half stand toward Ebal, half toward
Gerizim, with the priestly bearers and the Ark at the center of the covenant
assembly.  Women, children, resident aliens, elders, officers, and judges hear
blessing and curse (Josh 8:30--35).  The Ark is not said to stand on either
summit.  Nor should the passage's position be pressed into an exact road
chronology: the Old Greek and a Dead Sea Scroll witness place this unit at
different points in Joshua's sequence.

At Shiloh the Tent of Meeting is erected and the land is apportioned
(Josh 18:1--10).  Joshua does not name the Ark in that installation notice;
the opening of Samuel later places it there explicitly (1 Sam 3:3).  Judges
20:18, 26--28 preserves a further textual branch: Hebrew-oriented witnesses
place the Ark at Bethel, while the Clementine Vulgate and Douay tradition read
the house of God as Shiloh.  This study will not turn the two readings into two
successive stations.

The road has changed Israel's landscape: wilderness, river, walls, covenant
assembly, sanctuary.  Yet the deepest lesson has not changed.  The Ark goes
before the obedient, remains behind the presumptuous, and stands among the
people when the covenant is proclaimed.

\begin{studybox}{Presence is not presumption}
Compare Numbers 10:33--36 with Numbers 14:39--45.  In the first scene the Ark
goes before a people moving at God's command.  In the second it remains in camp
while zeal without obedience runs toward defeat.  Spiritual confidence is not
the certainty that God must bless whatever we undertake.  It is the freedom to
follow when he leads and to remain when he does not send.
\end{studybox}
